\section{Discussion and Limitations}
\label{sec:discussion}

The useful result is asymmetry, not aggregate ``GPU speedup.''  CPU/OpenSSL fits AES-GCM data writes, while GPU/cuPQC helps slack-tolerant ML-KEM maintenance.  AEGIS-Q trades peak background throughput for explicit publication, placement, QoS, and recovery boundaries; stale telemetry costs elastic progress, not the CPU/D/J/C foreground path.

The trust model remains kernel/FUSE-daemon trust plus backing-format checks.  The mount password is the root credential; privileged-local attacks, GPU side channels, and hardware-backed key release are outside the current claim.  TPM support is limited to replay-after-advance fail-closed behavior, not persistent PCR-bound lifecycle protection.

Deployment scope is intentionally narrow.  The frozen matrix covers plaintext, gocryptfs, dm-crypt, and AEGIS-Q; fscrypt is environment-blocked, so the matched baseline claim is measured mode alignment rather than exhaustive kernel comparison.  We evaluate read/write/\texttt{fsync}, scoped rename, directory \texttt{fsync}, and bounded writer races, but AEGIS-Q is not a general POSIX filesystem and excludes shared \texttt{mmap}/\texttt{msync}, full POSIX edge cases, broad application concurrency, or full cross-SoC portability.  The crash matrix covers selected daemon/app cut points and loopback lower-block faults, not physical power-loss or kernel-crash certification.

\subsection{Deployment takeaway}
The evaluated envelope is local: authenticated FUSE storage with CPU AES-GCM publication, slack-gated PQC maintenance, executor-local managed memory, and kernel/FUSE-daemon trust.  Figure~\ref{fig:first_page_qos} and Figure~\ref{fig:evaluation_summary}(b) reuse the same SQLite hero artifact; elastic background writes and append-log/cache-manifest remounts are the mounted behavior, not a separate deployment anecdote.  Jetson/CUDA/TPM dependencies limit portability, TPM evidence is fail-closed replay-after-advance rather than persistent PCR-bound key release, and thermal logs remain observations rather than energy-efficiency claims.  Future hardening is persistent rollback policy, wider QoS/workload comparison, and kernel-path acceleration only if authenticated publication and recovery oracles remain intact.
